\section{Conditional inference and scenario families}

\subsection{What the probability statement conditions on}

With the financial model, fitted policy and approximate values fixed, simulate independent outer paths. Let $X_i=G_\pi^{(i)}(M)/b$, with sample mean $\overline X$ and sample standard deviation $s_X$. A commonly used approximate upper limit is
\begin{equation}
 B_\pi=\overline X+t_{1-\alpha,n-1}\frac{s_X}{\sqrt n}.
\label{eq:pointwise}
\end{equation}
Since $\loss_\pi\leq\E[X_i]$, a valid upper limit for the expected gap is also an upper limit for policy loss. Student calibration is generally approximate for non-normal gaps; it is not a finite-sample guarantee. Finite variance, independent evaluation and sufficiently informative Monte Carlo samples remain necessary. More outer paths improve precision about the expected gap but do not remove its dual slack.

A conditional interval does not integrate over the uncertainty of policy training, financial calibration, contract conventions or sample selection. Training repetitions must be retained separately if their variation is investigated. Choosing a policy on the same evaluation outcomes used to advertise its bound can invalidate the selected-policy interpretation. Such a selection requires a separate evaluation design or an appropriate simultaneous guarantee.

\subsection{Coverage for a prescribed financial family}

For $m=|\Theta|$, compute a per-scenario limit at the maximal calendar using error allocation $\alpha/m$:
\begin{equation}
 B_\theta=\overline X_\theta+
 t_{1-\alpha/m,n_\theta-1}\frac{s_\theta}{\sqrt{n_\theta}}.
\label{eq:family}
\end{equation}
If each marginal limit has its stated coverage, the union bound and \eqref{eq:maximal} imply an event of probability at least $1-\alpha$ on which every prescribed scenario and smaller calendar is covered. With approximate marginal calibration, the joint statement inherits that approximation. Independence across scenarios is unnecessary for the union bound.

There is no extra multiplicity factor for calendars already dominated by the maximal endpoint. There is also no automatic extension to all policies, a continuous funding region or scenarios introduced after inspection. A missing required scenario prevents a complete-family conclusion; the family size must not shrink opportunistically. Financial quote compatibility is a different condition and must be reported separately from mathematical completeness of the conditional family.

\subsection{Pairing, independent runs and gap interactions}

Within a common-path experiment, I estimate policy gains or scenario effects from pathwise differences. For $D_i=X_i-Y_i$,
\begin{equation}
 \widehat{\operatorname{Var}}(\overline D)=\frac{s_D^2}{n}
 =\frac{s_X^2+s_Y^2-2s_{XY}}{n}.
\label{eq:paired}
\end{equation}
Discarding covariance would answer a different uncertainty question. Conversely, unrelated evaluation streams are not paired merely because policies or model labels coincide.

For independent mean estimates $R$ and $L$, possibly with different path counts,
\begin{equation}
 \widehat v=\frac{s_R^2}{n_R}+\frac{s_L^2}{n_L},\qquad
 \nu=\frac{\widehat v^2}
 {(s_R^2/n_R)^2/(n_R-1)+(s_L^2/n_L)^2/(n_L-1)}.
\label{eq:welch}
\end{equation}
The usual Welch interval uses $t_{1-\alpha/2,\nu}\sqrt{\widehat v}$. Degenerate zero-variance cases require explicit treatment. A difference between reported confidence endpoints is a descriptive comparison, not itself a confidence interval for a policy-value difference.

A financial/calendar interaction can be written in terms of expected quoted gaps as
\begin{equation}
 I_\theta=\{g_\theta(d_{\max})-g_0(d_{\max})\}
          -\{g_\theta(0)-g_0(0)\}.
\label{eq:interaction}
\end{equation}
When each financial effect is paired internally but the two calendar experiments are independent, add the variances of their paired mean effects. By \eqref{eq:decomposition}, $I_\theta$ combines interactions in true loss and dual slack. It does not identify either component alone and has no causal interpretation. Repeated dependent pointwise comparisons do not constitute a simultaneous test; intervals containing zero do not establish equivalence or exact additivity.
